\documentclass[instructions.tex]{subfiles}
\begin{document}
 
\chapter{Le respect humain est ordinairement la marque d'un esprit foible}
 
C'est une grande foiblesse d'être esclave du respect humain, parce que les discours du monde et ses jugemens sont pour l'ordinaire impuissans, inévitables, inconstans et trompeurs.

\primo Ils sont impuissans. Ce sont des paroles que le vent emporte, qui ne peuvent pas plus vous nuire que le vent qui court aux Indes. Ce qu'on dit de vous vous laisse toujours tel que vous êtes, pauvre ou riche, bon ou méchant; il ne peut vous rendre meilleur, ni plus mauvais. Vous êtes ce que vous êtes devant Dieu, et rien de plus. Ce n'est pas le monde, mais votre conscience que vous devez écouter en ce point.

Les discours du monde ne peuvent vous affliger qu'autant que vous le voulez. Qu'on vous critique et qu'on vous raille, vous pouvez souffrir et vous taire. On ne pourra jamais en dire de vous autant que vous pouvez en souffrir. David, tressaillant d'allégresse devant l'arche, se mettoit peu en peine des discours du peuple, et souffroit les railleries de sa femme. Laissez parler; en faisant votre devoir, vous serez en paix.

Si le monde attaque votre honneur, confiez-vous en Dieu; il sait le degré de réputation qui vous convient pour sa gloire; si vous lui êtes fidéle, il parlera pour vous; il prendra votre défense, comme il prit celle de Magdeleine contre le pharisien. Pourvu qu'il vous dise, comme à cette courageuse pénitente: Vos péchés sont remis, vous devez être tranquille. Conservez votre réputation par une vie sainte; c'est le moyen, dit saint Pierre, de faire taire l'impudence des hommes ignorans\footnote{1. Petr. 2.}. Si vous craignez les discours du monde, dit le Sage, vous succomberez\footnote{Qui timet hominem, cito corruet. Prov. 29.}. Mais, si vous craignez Dieu, rien ne sera capable de vous ébranler\footnote{Qui timet Dominum, nihil trepidabit. Eccli. 34.}.

\secundo Les jugemens et les discours du monde sont inévitables. Quoi que vous fassiez, vous serez désapprouvé de quelques-uns, et peut-être d'un grand nombre. Fussiez-vous le plus irréprochable, le plus saint des hommes, vous aurez des contradicteurs. Plus vous ferez votre devoir, plus la malignité s'élèvera contre vous; loin de vous émouvoir, vous devez vous y attendre. On a parlé contre Jésus-Christ, on a raillé sa conduite sainte: vous n'êtes pas meilleur que votre maître\footnote{Non est discipulus super magistrum. Luc. 6.}. Vous devez donc vous soucier peu des jugemens du monde. C'est une grande foiblesse, dit saint Martin de Brague, de craindre ce qu'on ne peut éviter et ce qui ne peut nuire.

\tertio Les jugemens du monde sont inconstans. Les hommes changent presque à tout moment; ils vous approuvent aujourd'hui, demain ils vous désapprouveront. Leur amitié et leur aversion, leur estime et leur indifférence, ont leur alternative. La même bouche souffle le chaud et le froid. Celui qui vous flatte aujourd'hui, demain vous flétrira. C'est donc une foiblesse de vous mettre en peine de ce qu'on dit de vous.

\quarto Enfin les jugemens des hommes sont souvent trompeurs et mal fondés. Combien de personnes vertueuses qu'on décrie, qui passent dans l'esprit d'un certain monde pour des gens méprisables, tandis que des scélérats sont honorés! Combien d'hommes savans et zélés passent chez quelques-uns pour des imprudens et de petits génies! Combien d'étourdis, de babillards ignorans, passent chez les autres pour de beaux esprits! Combien de femmes appliquées à sanctifier leurs familles, qu'on blâme et qu'on censure, tandis que des coquettes et des mondaines sont applaudies et préconisées! Jugez de là du peu d'estime qu'on doit faire des discours du monde, et combien il y a de foiblesse à les craindre.

Faites donc peu de cas du mépris des hommes et de leur approbation, et vous jouirez d'une paix inaltérable. Laissez-les penser, raisonner, murmurer. Si votre conscience est pure, que vous peuvent-ils? Peu m'importe, disoit saint Paul, ce que vous pensiez de moi; ma conscience ne me reproche rien, c'est Dieu seul qui doit me juger\footnote{Mihi autem pro minimo est ut a vobis judicer... Nihil mihi conscius sum... Qui autem judicat me, Dominus est. Cor. 4.}.

Concluons que le respect du monde ne doit pas être le mobile de nos actions, mais Dieu seul. Édifions le public, ne scandalisons personne, craignons Dieu, et laissons dire. Nous n'aurons jamais ni repos ni sainteté, si nous nous arrêtons à ce que le monde dira de nous, et si nous voulons régler notre conduite sur ses maximes et sur ses sentimens.


\section{Résolutions}

\primo Je profiterai des jugemens des hommes à mon égard, quand ils reprendront en moi un mal réel; alors je travaillerai à me corriger de ce défaut, de ce vice.

\secundo Mais quand ces jugemens tendront à me détourner de ce que je dois à Dieu, à ma conscience, à mon état, je n'en tiendrai pas plus de compte que du vent qui court aux Indes.

\prayer{Pardonnez-moi, ô mon Dieu, les péchés qu'un lâche respect humain m'a fait commettre dans le passé, et daignez m'affermir puissamment dans la résolution que je fais aujourd'hui de ne plus écouter ce honteux mobile.}

\end{document}
